\documentclass[conference]{IEEEtran}

\usepackage[utf8]{inputenc}
\usepackage[T1]{fontenc}
\usepackage{amsmath,amssymb}
\usepackage{booktabs}
\usepackage{tabularx}
\usepackage{cite}
\usepackage[hidelinks]{hyperref}

\newcommand{\SO}{\mathrm{SO}(3)}
\newcommand{\SE}{\mathrm{SE}(3)}

\begin{document}

\title{M2: Grasp Evaluation in MuJoCo\\[0.3em]
  {\normalsize Milestone report $\cdot$ ADLR WS 2026/27, project idea 4}}

\author{
  \IEEEauthorblockN{Georgii Fatakhov}
  \IEEEauthorblockA{Technical University of Munich}
  \and
  \IEEEauthorblockN{Awar Satar}
  \IEEEauthorblockA{Technical University of Munich}
}

\maketitle

\section{Goal and Setup}

M2 checks generated grasps physically in MuJoCo~\cite{todorov2012mujoco}. The completion criterion,
ground-truth grasps succeed at ${\gtrsim}\,90\,\%$, random grasps at ${\approx}\,0\,\%$, and
EquiGraspFlow~\cite{lim2024equigraspflow} is evaluated across all NFE, is met. The test cases are those of
the EMD evaluation in M1: 83 ACRONYM~\cite{eppner2021acronym} test objects in 3 rotations (249 cases),
full point clouds, 100 grasps per case.

\section{Approach}

\textbf{Gripper.} A floating Franka Panda hand (collision geometry from MuJoCo
Menagerie~\cite{zakka2022menagerie}) on a vertical, position-controlled lift axis, 40\,N per finger,
friction 1, with stiff contacts so that light objects are not squeezed through.

\textbf{Object.} A free body made of the convex parts of a CoACD~\cite{wei2022coacd} decomposition, which
keeps the cavities of mugs and bowls, with ACRONYM's mass.

\textbf{Protocol and metrics.} The open hand is placed at the grasp pose. A \emph{collision} is a
penetration of the object by more than 1\,mm; \emph{contact} means the object reaches into the volume
swept by the fingertip pads. Grasps with contact and no collision are simulated: close without gravity,
switch gravity on, lift by 15\,cm in 1\,s, hold for 1\,s. \emph{Success} means both fingers still touch
the object and it has risen by more than half the lift height; colliding grasps fail. Rates are fractions
of all grasps.

\textbf{Grasps} are random successful ACRONYM grasps (ground truth), samples of the model prior
(random), or EquiGraspFlow samples from the released checkpoint, with 1--20 Euler steps or the default
Runge--Kutta--Munthe-Kaas sampler (RK4-MK, 20 steps). As in M1, NFE counts evaluations of the guided
vector field (two network passes each). Latency is the median time to sample one grasp.

\textbf{Verification.} 38 unit tests on synthetic objects pass. Breaking the hand-frame mapping on
purpose drops ground-truth success on four objects to 10--43\,\%.

\begin{table}[!t]
  \centering
  \caption{Full test split, in \%. Latency of sampling one grasp in ms: NVIDIA L4 (same run) / RTX 5070 Ti
  (local subset run).}
  \label{tab:results}
  \begin{tabular}{@{}lrrrrr@{}}
    \toprule
    Grasps & NFE & Success & Collision & Contact & Latency \\
    \midrule
    Ground truth & -- & 98.9 & 0.0 & 100.0 & -- \\
    Random & -- & 0.1 & 22.1 & 5.6 & -- \\
    \midrule
    Euler & 1 & 0.7 & 56.0 & 20.2 & 97 / 37 \\
    Euler & 2 & 11.1 & 63.4 & 51.9 & 98 / 39 \\
    Euler & 5 & 62.0 & 28.2 & 91.1 & 115 / 53 \\
    Euler & 10 & 93.7 & 4.8 & 99.8 & 147 / 76 \\
    Euler & 20 & 96.2 & 2.5 & 99.9 & 212 / 120 \\
    RK4-MK & 80 & 96.4 & 2.5 & 100.0 & 579 / 411 \\
    \bottomrule
  \end{tabular}
\end{table}

\section{Results}

Ground truth succeeds at 98.9\,\%, random grasps at 0.1\,\%, mostly because the object is rarely between
the fingers (Table~\ref{tab:results}). EquiGraspFlow fails at NFE~1--2, where most grasps collide or miss,
reaches 93.7\,\% at NFE~10 and 96\,\% at NFE~20 and 80; the gap to ground truth is mostly Mug (90.0\,\% at
NFE~80). A local run on 5 objects per category agrees within 3 points.

\section{Limitations}

\textbf{The lift test is lenient.} ACRONYM-labelled failures pass at 62--92\,\% on four Mug and Pencil
objects (ground truth 90--100\,\%): objects weigh at most 0.42\,kg, the grip is strong, and the lift is
slow and without ACRONYM's shaking. Friction 0.3 drops them to 33--68\,\% (ground truth 93--100\,\%).
Success may thus saturate and understate differences between methods.

\textbf{Not comparable} with the paper's success rates (Isaac Gym, shaking gripper).

\textbf{Reproducibility.} The reference loader samples test point clouds with Open3D's unseeded generator:
two identical runs differed in 18 of 240 grasp outcomes, and seeding Open3D made them identical. This also
affects M1's EMD.

\section{Next Steps}

Seed Open3D, add a stricter test variant, evaluate M3 with this protocol.

\bibliographystyle{IEEEtran}
\bibliography{adlr}

\end{document}
